\documentclass[noauthor,noinstructornotes]{ximera}

\title{Week 11}
\author{Math 1193 Design Team}

\begin{document}
\begin{abstract}
    Activities for Week 11.
\end{abstract}
\maketitle

\begin{problem}
Let's play a similar game to the one we played two weeks ago! 
You'll pair up to think about this problem together, but play against me this time.
We'll start with a function, 
\(f\left(x\right)=2^x\). I'll select an \(x\)-value in secret and 
tell you the corresponding output. Your job is to guess what my 
\(x\)-value was. 

\begin{enumerate}
    \item If the output is 4, what is my \(x\)-value?
    \item What if my output is 2?
    \item 1?
    \item \(\frac{1}{2}\)?
    \item 64?
    \item \(\sqrt{2}\)?
\end{enumerate}

What process did you use to find these \(x\)-values? Be as specific as possible.

If you get bored, try to play my (the author's) role.
\end{problem}

\begin{problem}
Find a partner.

\begin{enumerate}
    \item Solve the following equations together.
    \begin{enumerate}
        \item \(3^x = 27\)
        \item \(4^x = 16\)
        \item \(4^x = 2\)
        \item \(5^x = \frac{1}{5}\)
        \item \(2^x = \frac{1}{16}\)
        \item \(9^x = \frac{1}{3}\)
    \end{enumerate}
    \item What do the equations have in common? What are the differences?
    \item Individually, write an equation of this type, then have your partner solve it.
\end{enumerate}
\end{problem}

\begin{problem}
We'll investigate powers of 5.

\begin{enumerate}
    \item Solve \(5^x = 25\). 
    \item Solve \(5^x = 125\).
    \item Solve \(5^x = 625\).
    \item Now we'll introduce a mystery number \(a\). The
    definition of \(a\) is that \(5^a = 50\). That is, \(a\)
    is the solution of the equation \(5^x = 50\).
    \begin{enumerate}
        \item Is \(a\) a whole number? (Whole numbers are 0, 1, 2, 3, 4, ...)
        \item What two consecutive integers does \(a\) fall between? (Some 
        examples of consecutive integers are \(-2\), \(-1\); or \(5\), \(6\); or \(101\), \(102\).)
    \end{enumerate}
\end{enumerate}

\begin{instructorNotes}
This really should hammer home that logarithms evaluated at positive numbers
do, in fact, produce real numbers. It tries to give students a place to 
put them on the number line. 
\end{instructorNotes}
\end{problem}

\end{document}
